\documentclass[12pt]{article}
\usepackage[utf8]{inputenc}
\usepackage[T1]{fontenc}
\usepackage{hyperref}
\title{Providencia Sign Language Language Resources for K-12 Curriculum Development}
\author{Compiled by Claude Sonnet 5.5 (claude-sonnet-5-5) for Andr\'as Kornai}
\date{}
\begin{document}
\maketitle

\section{Introduction}
Providencia Sign Language (Providence Island Sign Language, ISO 639-3 \texttt{prz}) is a shared sign language of Providencia, in the San Andr\'es, Providencia y Santa Catalina department of Colombia. Sources describe it as used by a small number of deaf islanders (about 20 per one project) and by many hearing islanders. Sources disagree on dates and numbers. Recent research asks whether it is better seen as emergence or revitalization within Colombian Sign Language (LSC) contexts, and the signing has little prestige in the multilingual island setting.

\section{Primary Educational Resources}
No curricula, textbooks or school programmes in Providencia Sign Language were found in the searches.

\section{Linguistic and Reference Materials}
\begin{itemize}
\item Revisiting the shared sign language of Providence Island (ELDP-funded documentation project): \url{https://cultureincrisis.org/projects/revisiting-the-shared-sign-language-of-providence-island}
\item Article in \emph{Forma y Funci\'on} (Universidad Nacional de Colombia) on the language and its situation: \url{https://revistas.unal.edu.co/index.php/formayfuncion/article/view/101071}
\item Same article, PDF: \url{https://revistas.unal.edu.co/index.php/formayfuncion/article/download/101071/89434/655006}
\item Sampling signers on Providence Island: reflections on a small-scale documentation project: \url{https://redcol.minciencias.gov.co/Record/UCART2_6ba0cc445dbbbbd6f73826825631b0da/Details}
\item Braithwaite (2017), SCL Occasional Paper 38: \url{https://scl-online.net/publications/SCL_OP_38_Braithwaite_2017.pdf}
\item Repository record at Universidad Nacional de Colombia: \url{https://repositorio.unal.edu.co/handle/unal/56239}
\item Verbix entry: \url{https://docs.verbix.com/EndangeredLanguages/ProvidenciaSignLanguage}
\item Rosetta Project entry: \url{https://rosettapanglossia.longnow.org/index.php?title=Prz}
\end{itemize}

\section{Religious and Community Resources}
None found.

\section{Digital and Audio Resources}
No public video corpus or dictionary was found in the search results.

\section{Curriculum Development Context}
The language community is very small and the documentation is mostly academic. A curriculum would need community-led video documentation, work with deaf signers and their families, and coordination with Colombian Sign Language education on the island.
\end{document}